\documentclass[a4paper]{article}
\usepackage{graphicx}
\usepackage{amsmath,amssymb}
\usepackage{hyperref}
\usepackage{minted}
\usepackage{multirow}
\usepackage{float}
\usepackage{todonotes}
\usepackage{forest}
\usepackage{xstring}
\input{/home/donkarlo/Dropbox/repo/nd_language_ai_project/src/nd_language_ai/formal/computer/programming/tex/latex/code_clips/external_dot.tex}
\title{}
\date{}

\begin{document}
    \maketitle
    
    
    \section{Bewertungskriterien „Sprechen“}
        
        \textbf{Deutsch:}
        Im Subtest „Sprechen“ wird die mündliche Sprachkompetenz der Teilnehmer und Teilnehmerinnen beurteilt. Dabei wird zwischen aufgabenbezogenen und sprachbezogenen Kriterien unterschieden. Die Kriterien basieren auf dem \textit{Gemeinsamen europäischen Referenzrahmen für Sprachen} und auf \textit{Profile Deutsch}. Folgende Kriterien werden berücksichtigt:
        
        \textbf{English:}
        In the subtest “Speaking”, the oral language competence of the participants is assessed. A distinction is made between task-related criteria and language-related criteria. The criteria are based on the \textit{Common European Framework of Reference for Languages} and on \textit{Profile Deutsch}. The following criteria are taken into account:
        
        \subsection{Inhaltliche Angemessenheit / Content Appropriateness}
            
            I Aufgabenbewältigung (Task accomplishment)
        
        \subsection{Sprachliche Angemessenheit / Linguistic Appropriateness}
            
            II Aussprache / Intonation (Pronunciation / Intonation)
            III Flüssigkeit (Fluency)
            IV Korrektheit (Correctness)
            V Wortschatz (Vocabulary)
            
            \textbf{Deutsch:}
            Der Subtest „Sprechen“ besteht aus fünf Teilaufgaben. Bei jeder Teilaufgabe wird bewertet, inwiefern das Kriterium der Aufgabenbewältigung erfüllt wurde. Dagegen beziehen sich die Kriterien II–V auf die gesamte mündliche Leistung, die die Teilnehmer und Teilnehmerinnen während des Prüfungsgesprächs (Teil 1–3) erbringen.
            
            \textbf{English:}
            The subtest “Speaking” consists of five subtasks. In each subtask, the extent to which the task accomplishment criterion is fulfilled is assessed. Criteria II–V refer to the overall oral performance demonstrated by the participants during the exam conversation (parts 1–3).
        
        \subsection{Inhaltliche Angemessenheit / Content Appropriateness}
            
            \begin{table}[H]
                \centering
                \begin{tabular}{|p{3cm}|p{4cm}|p{4cm}|p{4cm}|}
                    \hline
                    & B1 & A2 & A1 \\
                    \hline
                    
                    Teil 1A
                    
                    \textbf{Deutsch:} Kann sich vorstellen und dabei auch detailliertere Informationen vortragen.
                    
                    \textbf{English:} Can introduce themselves and provide more detailed information.
                    &
                    \textbf{Deutsch:} Kann sich vorstellen und dabei knappere, allgemeine Informationen geben.
                    
                    \textbf{English:} Can introduce themselves and give shorter, general information.
                    &
                    \textbf{Deutsch:} Kann sich vorstellen und dabei Informationen unverbunden vortragen.
                    
                    \textbf{English:} Can introduce themselves and present information in a simple, unconnected way.
                    \\
                    \hline
                    
                    Teil 1B
                    
                    \textbf{Deutsch:} Kann auf Nachfragen relativ spontan und ausführlich antworten.
                    
                    \textbf{English:} Can respond to follow-up questions relatively spontaneously and in detail.
                    &
                    \textbf{Deutsch:} Kann auf Nachfragen knapp und/oder nur teilweise verständlich antworten.
                    
                    \textbf{English:} Can answer follow-up questions briefly and/or only partially understandably.
                    &
                    \textbf{Deutsch:} Kann auf Nachfragen, die langsam, deutlich und in direkter Sprache gestellt sind, mit einzelnen Wörtern oder auswendig gelernten Wendungen antworten.
                    
                    \textbf{English:} Can answer follow-up questions asked slowly and clearly using single words or memorized phrases.
                    \\
                    \hline
                    
                    Teil 2A
                    
                    \textbf{Deutsch:} Kann die Hauptinhalte eines Fotos und auch Einzelheiten benennen.
                    
                    \textbf{English:} Can name the main contents of a photo and also details.
                    &
                    \textbf{Deutsch:} Kann die Hauptinhalte eines Fotos knapp und sehr allgemein benennen.
                    
                    \textbf{English:} Can name the main contents of a photo briefly and very generally.
                    &
                    \textbf{Deutsch:} Kann die Hauptinhalte eines Fotos in sehr wenigen Worten andeuten.
                    
                    \textbf{English:} Can indicate the main contents of a photo in very few words.
                    \\
                    \hline
                    
                    Teil 2B
                    
                    \textbf{Deutsch:} Kann auf Nachfrage eigene Erfahrungen teilweise detailliert berichten.
                    
                    \textbf{English:} Can describe personal experiences in some detail when asked.
                    &
                    \textbf{Deutsch:} Kann auf Nachfrage eigene Erfahrungen knapp und allgemein berichten.
                    
                    \textbf{English:} Can briefly describe personal experiences when asked.
                    &
                    \textbf{Deutsch:} Kann auf Nachfrage mit einzelnen Wörtern und sehr knappen Äußerungen antworten.
                    
                    \textbf{English:} Can respond with single words or very short utterances when asked.
                    \\
                    \hline
                    
                    Teil 3
                    
                    \textbf{Deutsch:} Kann ein Gespräch beginnen und in Gang halten. Kann im Gespräch spontan etwas planen, Ideen und Meinungen mitteilen und Vorschläge machen sowie darauf reagieren.
                    
                    \textbf{English:} Can start and maintain a conversation, plan something spontaneously in conversation, express ideas and opinions, make suggestions and respond to them.
                    &
                    \textbf{Deutsch:} Kann Fragen stellen und beantworten, versteht aber kaum genug, um selbst das Gespräch in Gang zu halten.
                    
                    \textbf{English:} Can ask and answer questions but usually cannot maintain the conversation independently.
                    &
                    \textbf{Deutsch:} Ist darauf angewiesen, dass Dinge langsamer wiederholt oder umformuliert werden.
                    
                    \textbf{English:} Depends on others repeating or reformulating things more slowly.
                    \\
                    \hline
                \end{tabular}
            \end{table}
            
            
            \newpage
        
        \subsection{Sprachliche Angemessenheit / Linguistic Appropriateness}
            
            \begin{table}[H]
                \centering
                \begin{tabular}{|p{3cm}|p{4cm}|p{4cm}|p{4cm}|}
                    \hline
                    & B1 & A2 & A1 \\
                    \hline
                    
                    II Aussprache / Intonation
                    
                    \textbf{Deutsch:} Spricht gut verständlich, auch wenn ein fremdsprachiger Akzent teilweise offensichtlich ist.
                    
                    \textbf{English:} Speaks clearly and understandably, even if a foreign accent is sometimes noticeable.
                    &
                    \textbf{Deutsch:} Spricht im Allgemeinen klar genug, auch wenn der Gesprächspartner gelegentlich um Wiederholung bitten muss.
                    
                    \textbf{English:} Generally speaks clearly enough, though the listener may occasionally ask for repetition.
                    &
                    \textbf{Deutsch:} Aussprache aus begrenztem Repertoire; kann von Muttersprachlern verstanden werden.
                    
                    \textbf{English:} Pronunciation from a limited repertoire; can be understood by native speakers.
                    \\
                    \hline
                    
                    III Flüssigkeit
                    
                    \textbf{Deutsch:} Kann sich ohne viel Stocken verständlich ausdrücken.
                    
                    \textbf{English:} Can express themselves understandably without much hesitation.
                    &
                    \textbf{Deutsch:} Kann vertraute Themen bewältigen, obwohl er/sie häufig stockt.
                    
                    \textbf{English:} Can manage familiar topics but often hesitates.
                    &
                    \textbf{Deutsch:} Kann sehr kurze, isolierte Äußerungen produzieren.
                    
                    \textbf{English:} Can produce very short, isolated utterances.
                    \\
                    \hline
                    
                    IV Korrektheit
                    
                    \textbf{Deutsch:} Gute Beherrschung der grammatischen Strukturen trotz Einflüssen der Muttersprache.
                    
                    \textbf{English:} Good command of grammatical structures despite influence from the first language.
                    &
                    \textbf{Deutsch:} Macht noch systematische Fehler, aber die Aussage bleibt klar.
                    
                    \textbf{English:} Still makes systematic errors but meaning remains clear.
                    &
                    \textbf{Deutsch:} Beherrscht einige wenige einfache Strukturen auswendig gelernt.
                    
                    \textbf{English:} Uses a few simple memorized grammatical structures.
                    \\
                    \hline
                    
                    V Wortschatz
                    
                    \textbf{Deutsch:} Verfügt über einen ausreichend großen Wortschatz, um die meisten Alltagsthemen zu äußern.
                    
                    \textbf{English:} Has a sufficiently large vocabulary to discuss most everyday topics.
                    &
                    \textbf{Deutsch:} Verfügt über genügend Wortschatz für einfache Grundbedürfnisse.
                    
                    \textbf{English:} Has enough vocabulary for basic everyday needs.
                    &
                    \textbf{Deutsch:} Beherrscht einzelne Wörter und kurze Sätze für konkrete Alltagssituationen.
                    
                    \textbf{English:} Uses individual words and short sentences for concrete everyday situations.
                    \\
                    \hline
                
                \end{tabular}
            \end{table}
    % \bibliographystyle{plain}
    % \bibliography{/home/donkarlo/Dropbox/repo/shared_research/bibliography/bibliography.bib}
\end{document}